\introChapter{พีชคณิตบูลีน (Boolean Algebra)}

\begin{description}[leftmargin=2.5cm, style=nextline]
    \item[รหัสวิชา] 4091606
    \item[ชื่อวิชา] คณิตศาสตร์สำหรับคอมพิวเตอร์
    \item[หน่วยกิต] 3 (3-0-6)
    \item[บทที่] 6
    \item[ชื่อบท] พีชคณิตบูลีน (Boolean Algebra)
    \item[จำนวนชั่วโมง] 9 ชั่วโมง
\end{description}

\vspace{1em}
\noindent วัตถุประสงค์การเรียนรู้ประจำบท\\[0.3em]
\noindent เมื่อนักศึกษาเรียนบทนี้จบแล้ว นักศึกษาจะสามารถ
\begin{enumerate}[topsep=0.3em, itemsep=0.25em]
    \item อธิบายสัจพจน์และกฎของพีชคณิตบูลีนได้
    \item ลดรูปสมการบูลีนโดยใช้กฎพีชคณิตบูลีนได้
    \item สร้างและวิเคราะห์วงจรประตูตรรกะจากนิพจน์บูลีนได้
    \item ใช้แผนที่การ์โนในการลดรูปฟังก์ชันบูลีนได้
    \item เชื่อมโยงพีชคณิตบูลีนกับการออกแบบวงจรดิจิทัลและการเขียนโปรแกรมได้
\end{enumerate}

\vspace{1em}
\noindent หัวข้อที่สอน\\[0.3em]
\begin{enumerate}[topsep=0.3em, itemsep=0.25em]
    \item สัจพจน์และกฎของพีชคณิตบูลีน
    \item การลดรูปนิพจน์บูลีน
    \item วงจรประตูตรรกะ
    \item แผนที่การ์โน
    \item การประยุกต์ใช้พีชคณิตบูลีนในวงจรดิจิทัล
\end{enumerate}

\vspace{1em}
\noindent สื่อการสอน
\begin{enumerate}[topsep=0.3em, itemsep=0.25em]
    \item สไลด์ประกอบการสอน
    \item ซอฟต์แวร์จำลองวงจรลอจิก
    \item แบบฝึกหัดท้ายบท
\end{enumerate}
